\documentclass{situationist-preprint}

\title{Compatibilism Without Determinism}
\msid{ACT-01-d}
\msfamily{ACT-01}
\msdiscipline{Action, Agency and Free Will}
\mstype{Consequence}
\mskeywords{compatibilism; determinism; indeterminism; reasons-responsiveness; reliability; Hobart}

\begin{document}
\maketitle

\begin{abstract}
Compatibilism is usually presented as a view about determinism: free will would survive if determinism were true. A tradition running from Hume to Hobart went further and held that free will requires determination, so that compatibilism stands or falls with determinism. We derive from the compatibility thesis of ACT-01 that this further claim is false. Compatibilism, understood as the view that free will consists in owned reasons-responsive control, does not require determinism, does not need to take a stand on it, and is not refuted by the discovery that the laws are indeterministic. What free will requires is reliability in the mechanism of practical reasoning, and reliability is a property that deterministic and indeterministic processes can equally have.
\end{abstract}

\section{From compatibility to independence}

ACT-01-a defines free action as action from a mechanism that is the agent's own and reasons-responsive (Definitions~2.2--2.4), and proves from an independence assumption (Assumption~3.1) that free will is compatible with determinism and with indeterminism (Theorem~3.2). Its Corollary~3.3 states that whether we have free will does not depend on whether determinism is true. We draw out three consequences for compatibilism itself.

\begin{definition}\label{d:comp}
\emph{Compatibilism} is the thesis that free will is compatible with determinism. \emph{Determinist compatibilism} is the further thesis that free will requires determinism: that no agent in an indeterministic world acts freely, or that she does so only in so far as her action is determined.
\end{definition}

Determinist compatibilism was the classical form of the view. \citet[\S VIII]{Hume1748} argued that liberty, as opposed to necessity, would amount to chance, and that actions not necessarily connected with the agent's character could not redound to her credit or blame. \citet{Hobart1934} gave the thesis its title: free will involves determination and is inconceivable without it.

\begin{proposition}\label{p:notreq}
Free will does not require determinism. Determinist compatibilism is false.
\end{proposition}

\begin{proof}
By Theorem~3.2 of ACT-01-a, free will is compatible with indeterminism: there is an indeterministic world containing an agent who has free will. So it is not the case that every free agent inhabits a deterministic world.
\end{proof}

The proof uses the second half of the independence assumption, which ACT-01-a supports by the observation that macroscopic processes of deliberation can be reliable although the microphysical processes that realize them are not deterministic. The weaker reading of Definition~\ref{d:comp}, on which an action is free only in so far as it is determined, fails for the same reason: what reasons-responsiveness requires is that the mechanism would act on sufficient reasons in a range of situations, and \enquote{would} here is satisfied by a mechanism that does so with high reliability rather than with nomic necessity.

\begin{proposition}\label{p:neutral}
The case for compatibilism need not take a stand on whether determinism is true.
\end{proposition}

\begin{proof}
By Corollary~3.3 of ACT-01-a, whether we have free will does not depend on whether determinism is true. The case that we have it therefore rests on facts, the presence in normal human agents of owned, reasons-responsive mechanisms, that are established without settling the question of determinism.
\end{proof}

\begin{proposition}\label{p:indetfree}
If the fundamental laws are indeterministic, compatibilism is not thereby refuted, and its account of free action applies unchanged.
\end{proposition}

\begin{proof}
Compatibilism, as defined, is a modal thesis about the compatibility of free will with determinism; the falsity of determinism leaves it untouched. The account of free action on which ACT-01-a bases it mentions neither determinism nor indeterminism, so it applies to an indeterministic world without revision.
\end{proof}

\section{What remains of Hobart's thought}

Hume's and Hobart's insight survives in a corrected form. They saw that an action that varied independently of the agent's reasons and character would not be hers to answer for. The correct requirement is reliability, not determination: the action must vary with the agent's reasons, in the pattern described in Definition~2.2 of ACT-01-a, and this pattern can be realized by an indeterministic process whose chances are concentrated on the options the reasons favor. Indeterminism threatens free will only where it intervenes between reasons and action so as to make the action vary independently of them, which is the luck objection that ACT-15 answers for the case in which the agent tries to bring about each of the competing outcomes.

\begin{objection}[Libertarian appropriation]
If free will is compatible with indeterminism, then compatibilism collapses into a weak libertarianism, and the distinction between the views disappears.
\end{objection}

\begin{reply}
Libertarianism holds that free will requires indeterminism \citep{Kane1996}. The first half of Theorem~3.2 of ACT-01-a implies that it does not, and Proposition~\ref{p:notreq} that free will does not require determinism either: it requires neither. The resulting view is not a compromise between the two traditional positions; it denies the shared premise of the determinist compatibilist and the libertarian, that the question of determinism bears on whether we are free. \citet{List2019} and \citet{Vihvelin2013} defend views with the same structure.
\end{reply}

\section{Conclusion}

From the compatibility thesis of ACT-01 it follows that compatibilism neither requires determinism nor needs to know whether it is true. The classical view that freedom requires determination confused reliability with necessitation. What free will requires is a mechanism that reliably tracks reasons and is the agent's own, and the laws of nature, deterministic or not, leave room for it.

\begin{thebibliography}{9}
\bibitem[Hobart(1934)]{Hobart1934} Hobart, R.~E. (1934). Free will as involving determination and inconceivable without it. \emph{Mind}, 43(169), 1--27.
\bibitem[Hume(1748)]{Hume1748} Hume, D. (1748). \emph{An Enquiry Concerning Human Understanding}, Section VIII.
\bibitem[Kane(1996)]{Kane1996} Kane, R. (1996). \emph{The Significance of Free Will}. Oxford University Press.
\bibitem[List(2019)]{List2019} List, C. (2019). \emph{Why Free Will Is Real}. Harvard University Press.
\bibitem[Vihvelin(2013)]{Vihvelin2013} Vihvelin, K. (2013). \emph{Causes, Laws, and Free Will: Why Determinism Doesn't Matter}. Oxford University Press.
\end{thebibliography}

\end{document}
